
%
%
%


\section {Introduction}\label{Intro}

\mn


For a set $T$ we use $\Omega(T)$ for the discrete cube
$\{0,1\}^T$ and $\mu_T$ for the uniform probability measure
on $\Omega (T)$.
In this paper
$f$ will always be a Boolean function on $\gO([n])$
(that is, $f:\gO([n]) \to \{0,1\}$, where, as usual,
$[n]=\{1\dots n\}$).
We write $\mu$ for $\mu_{[n]}$ and
$\mu(f)= \mu (\{ x:f(x)= 1 \})$).
We reserve $x,y$ for elements of $\Omega([n])$
and set $|x| = \sum x_i$.





Following Ben-Or and Linial \cite {bl} we define,
for a given $f$ and
$S \subset [n]$, the \emph{influence of $S$ toward one}
to be
\beq{I+}
I^+_S(f) = \mu_{[n]\sm S}( \{u \in \Omega ([n]\backslash S):
\exists v \in \Omega(S), f(u,v)=1\})- \mu(f) .
\enq
Similarly, the influence of $S$ toward zero is
\beq{I-}
I^-_S(f) = \mu_{[n]\sm S}( \{u \in \Omega ([n]\backslash S):
\exists v \in \Omega(S), f(u,v)=0\})- (1-\mu(f) )
\enq
and the (total) influence of $S$ is
$$I_S(f)=I_S^+(f)+I_S^-(f).$$


Suppose $\mu(f) =1/2$.
It then follows from a theorem of Kahn, Kalai and Linial
\cite{kkl}
(\expref{Theorem}{t:kkl} below, henceforth ``KKL'')
that for every $a\in (0,1) $ there is an $S \subset [n]$ of size
$an$ with
$I^+_S(f)\ge 1/2-n^{-c}$, where $c>0$ depends on $a$.
(See \expref{Theorem}{t:ra1}.)
Benny Chor conjectured in 1989 \cite {benny} that one can in fact
achieve
$I^+_S(f)\ge 1/2-c^n$ (where, again, $c<1$ depends on $a$).
The conjecture has been ``in the air'' since that time,
though as far as we know
it has appeared in print only in
\cite{k:nati,Nati}.

%
In this paper we disprove
Chor's conjecture and another, similar conjecture
from the same period. In the subsequent part II,
we improve the preceding consequence of KKL.

For our purposes the subtracted terms in
equations %
%
\eqref{I+} and
\eqref{I-} are mostly a distraction, and it sometimes seems
clearer to speak of $J^+_S(f):=I^+_S(f)+\mu(f) $
and $J^-_S(f):=I^-_S(f)+(1-\mu(f) )$.
Thus, for example, $J^+_S(f)$ is the probability that
a uniform setting of the variables in $[n]\sm S$
doesn't force $f=0$, and Chor's conjecture predicts
an $S$ with $J^+_S(f)\ge 1-c^n$.
The following statement shows that this need not be the case.

%
\begin{thm}\label{prop1}
For any fixed $\ga,\gd \in (0,1)$
%
there exists $C\in\R$ such that for
all   %
large enough $n$
%
there exists 
%
%
%
%
an    %
$f$ such that  
$\mu(f) =\ga$ and $J_S^+(f)<1-n^{-C}$
for every $S\sub [n]$ of size %
%
at most $(1/2-\gd)n\,.$ 
%
%
\end{thm}

\nin
We should note that one cannot expect to go much
beyond $|S|=(1/2-\gd)n$; for example if $\mu(f)  =1/2$,
then it follows from the ``Sauer--Shelah Theorem''
(\expref{Theorem}{t:sas}) that there is
an $S$ of size $n/2$ with
$J_S^+(f) = 1$.

Another consequence of KKL (see \expref{Theorem}{t:ra2} below)
is that there is a $\gb>0$ such that for any $f$ with $\mu(f) > n^{-\beta}$
there is an $S$ of size (say) $0.1 n$ with influence $1-o(1)$.
A conjecture of the third author, again from
the late 80s, asserts
that the same conclusion holds even assuming only
$\mu(f) > (1-\eps)^n$
for sufficiently small $\eps$.
This conjecture turns out to be false as well.  %


%
\begin{thm}\label{prop2}
For any fixed $\eps, \gd>0$
%
there exists $\ga > 0$ such that for all large enough $n$
%
there exists an $f$ %
%
such that
$\mu(f) > (1-\eps)^n$ and no set of %
%
size at most $(1/2-\gd)n$   %
%
has influence
%
toward   %
1 more than $\exp[-n^\ga]$.
\end{thm}


\mn
This can be strengthened a bit to require
$\mu(f)  > \exp[-n^{1-c}]$ for some fixed $c=c_\gd>0$.





The examples proving \expref{Theorems}{prop1} \expref{and}{prop2} 
are given in \expref{Section}{examples}.
Each of these is of the form %
$f= \bigwedge_{i=1}^m C_i$,  %
where the %
$C_i$  %
are random %
disjunctions %
of $k$ literals
using $k$ distinct variables (henceforth ``$k$-clauses'').
These $f$'s, %
%
which may be thought of as
variants of the ``tribes''
construction of Ben-Or and Linial
(see below),
were inspired by a paper of
Ajtai and Linial \cite {AL} and share with it the following
curious feature.
%
It is  %
easy to see that any $f$ can be converted to a
\emph{monotone}
(\ie, increasing) $f'$ with $\mathbb E (f')=\mu(f) $ and each
influence ($I_S^+$ and so on) for $f'$ no larger than
the corresponding influence for $f$;
thus %
it is  %
natural to look for $f$'s %
%
with small influences
among the increasing functions.
But the present random examples, like those of \cite{AL},
do not do this, and %
it is %
not easy to see what one gets by monotonizing them.

\subsection *{Improving the consequences of KKL}

While the preceding, rather optimistic conjectures turn out to be false, 
in part II  we show
that the first of the aforementioned consequences of KKL can be 
improved. (The results and proofs of both parts were given together in
the arXiv publication \cite {BKK21}.) 
%
%

\noindent
\begin{thm}\label{prop3}

For each $C>0$, there is a $\delta>0$ such that for any $f$ with 
$\mu(f) >n^{-C}$, there
is an $S\subset [n]$ of size at most $(1/2-\delta)n$ with $J_{S}^{+}(f)>.9.$
\end {thm}
\nin
(Of course, as elsewhere in this discussion, ``.9'' could be any preset $\rho <1$.)

Note that the gap between \expref{Theorems}{prop2} \expref{and}{prop3} is substantial 
and our modest progress is likely not the final word on the problem.
For example, could it be that there is some fixed $\gb$ such that there
are $f$'s with %
%
%
$\mu(f) > n^{-\gb}$ for which no $S$ of size $0.1 n$ has influence $\gO(1)$?
We will discuss this question further in the next section.


%
%
%
The Proof of Theorem 1.3 (\cite {BKK21}) goes roughly as follows.
We employ two strategies, both variants of the analysis in \cite{kkl}.
The first 
uses the total influence, assumed sufficiently large.
If at some point this total influence becomes ``small,'' 
we switch to a different procedure that
combines the incremental argument from \cite{kkl} with the 
Sauer--Shelah %
theorem.  %

\section {Background and perspective}
\label{s:bac}

\subsubsection *{Influence}

We write $I_\ell (f)$ for $I_{\{\ell\}}(f)$.  A form of
the classic edge isoperimetric inequality for Boolean
functions is  
%

\begin{thm}\label{Tiso}
For any (Boolean function) $f$ %
%
 with $\mu(f) =t$,
\begin {equation}
\label {e:harper}
I(f) := \sum_{k=1}^n  I_\ell (f) \ge  2t \log_2 (1/t).
\end {equation}
\end {thm}

\mn
(This convenient version is easily derived from the precise
statement, due to Hart \cite{Hart}; see also \cite[Sec. 7]{KK07}
for a simple inductive proof.)

\mn
While \eqref{e:harper} is exact or close to exact
(depending on $t$), it typically gives
only a weak lower bound on the maximum of the %
$I_\ell (f)$'s,   %
namely
\begin {equation}
\label {e:harper-max}
\max_\ell I_\ell (f) \ge  2t \log_2 (1/t)/n.
\end {equation}
For $t$ not too close to 0 or 1, the
following
statement from \cite{kkl} gives
better information.


\begin{thm} [KKL]
\label {t:kkl}
There is a fixed $c>0$ such that
for any
$f$ with $\mu(f) =t$,
there is an $\ell \in [n]$ with
\begin {equation}
\label {e:kkl}
I_\ell (f) \ge c t (1-t)\log n/n.
\end {equation}
\end {thm}
\nin
Recall that $J_\ell (f)=\mu(f) +I_\ell (f)$. Repeated application of \expref{Theorem}{t:kkl} gives
the following two corollaries.

\begin {thm}
\label {t:ra1}
For all $a, t\in (0,1)$ there is a c such that
for any
$f$ with $\mu(f) =t$ there
is an $S\sub [n]$ with $|S|\leq an$ and
$$J^+_S(f) \ge 1 - n^{-c}$$ (that is,
$I^+_S(f) \ge (1-t) - n^{-c}$).

\end {thm}
\nin
Similarly (either by the same argument or by applying
\expref{Theorem}{t:ra1} to the function $1-f(x)$)
there is a small $S'$
with
$J^-_{S'}(f) \ge 1 - n^{-c}$ (\ie, $I^-_{S'}(f) \ge t - n^{-c}$),
and combining these observations we find that there is
in fact a small $S''$ (\eg, $S \cup S'$) with
$I_{S''}(f) \ge 1-n^{-c}$.


\begin {thm}
\label {t:ra2}
For every $\delta, \epsilon >0$,
there is an $\ga>0$ such that for large enough
$n$ and any
$f$ with $\mu(f) \geq  n^{-\ga}$,
there is an $S\sub [n]$ with $|S|=\delta n$ and
 $$J^+_S(f) \ge 1-\epsilon.$$
\end {thm}


\mn
The conjecture of
Chor stated in \expref{Section}{Intro}
asserts that the $n^{-c}$ in
\expref{Theorem}{t:ra1} can be replaced
by something exponential in $n$,
and the conjecture stated before
\expref{Theorem}{prop2} proposes a similar weakening
of the $n^{-\ga}$ lower bound on $\mu(f) $ in
\expref{Theorem}{t:ra2}.
As already noted, we
will show below that these conjectures are incorrect.



\subsubsection *{Tribes}
The original ``tribes'' examples of Ben-Or and Linial \cite{bl}
are Boolean functions of the form
$f= \bigvee_{i=1}^m C_i$, where the   %
``tribes'' $C_i$ are %
conjunctions %
of $k$ (distinct) variables
and
each variable belongs to exactly one tribe.
The dual of such an $f$ (so ``dual tribes'') is
%
$g= \bigwedge_{i=1}^m D_i$, where  %
$D_i$ is the %
disjunction %
of the variables in $C_i$
(so again, each variable belongs to exactly one $D_i$).

When $k=\log n - \log \log n - \log \ln (1/t)$,
we have $1-\mu(f)  = \mathbb E (g) \approx t$
(where $\log = \log_2$ and $f,g$ are as above).
For fixed $t\in (0,1)$ both constructions
show that \expref{Theorem}{t:kkl}
is sharp (up to the value of $c$).


On the other hand,
when $t=O(n^{-c})$ for a fixed
$c>0$, $f$ shows that \eqref{e:harper-max}
is tight up to a multiplicative constant, depending on $c$;
for example,
$k=2\log n-\log \log n$ gives
$\mu(f)  \approx 1/(2n)$ and
$I_\ell(f)\approx 2\log n/n^2=\Theta(\mu(f) \log(1/\mu(f) )/n)$
for each $\ell$.
(In contrast, for $\mathbb E (g)\approx 1/n$, we should
take $k= \log n-2\log \log n-1$, in which case
$I_\ell (g)=\Theta (\log^2n/n^2)$ and \eqref{e:harper-max}
is off by a log.)

For $f$ (again, as above) with $\mu(f)  \in (\gO(1),1-\gO(1))$,
there are sets of size $\log n$ with large influence
%
toward 1,  %
while no set of size $o(n/\log n)$ has
influence $\gO(1)$ %
toward 0.
(The corresponding statement with
the roles of 0 and 1 reversed holds for $g$.)
The Ajtai-Linial construction
mentioned in the introduction shows that there are Boolean functions
$h$ with $\mathbb E (h) \approx 1/2$ and $I_S(h) < o(1)$ for
every $S$ of size $o(n/\log^2n)$.



\subsubsection *{Trace}


We now briefly consider influences from a different point of view.
For a set $X$ let
$2^X = \{S: S \subset X\}$,
$\binom{X}{k}=\{S \subset X:|S|=k\}$,
$\binom{X}{< k}=\{S \subset X:|S| < k\}$ and 
$\binom{n}{<k}=\sum_{i=0}^{k-1} \binom{n}{i}.$
For $\f\subset 2^X$
and $Y \subset X$,
the \emph{trace} of $\cal F$ on $Y$ is
$${\cal F}_{|Y} = \{S \cap Y:
S \in {\cal F}\}.$$

Let $X = [n]$.
%
The ``Sauer--Shelah Theorem'' (below)
determines, for every $n$ and $m$, the minimum $T$
such that for each $\f\sub 2^X$ of size $T$ there is some
$Y \in \binom{X}{m}$ on which the trace of $\f$ is \emph{complete,}
meaning ${\cal F}_{|Y}=2^Y$. Such a $Y$
is said to be \emph{shattered} by ${\cal F}$.

\begin {thm}[The Sauer--Shelah Theorem]
\label {t:sas}
If ${\cal F} \subset 2^{X}$ and
$|\f| > \binom{n}{< r} $,
%
then $\f$ shatters some $Y\in \binom{X}{r}$.
\end {thm}

\nin
That this is sharp is shown by
${\cal F} = \binom{X}{< r}$, the \emph{Hamming ball}
of radius $r-1$ about
$\0$ with respect to the usual Hamming metric on
$2^{X}\equiv \gO(X)$.

\expref{Theorem}{t:sas} was proved
around the same time by Sauer \cite {Sa},
Shelah and Perles \cite{Sh}, and Vapnik and Chervonenkis \cite {VC}.
%
It has many connections, applications and extensions
in combinatorics, probability theory, model theory, analysis, statistics
and other areas.
%

We identify $\gO([n])$ and $2^{[n]}$ in the usual way.
The connection between traces and influences is as follows.
Let $f$ be a Boolean function on $\gO([n])$
and
${\cal F} =f^{-1}(1)$.
It is easy to see
that for
$S\sub [n]$ and $T = [n] \backslash S$,
$$J^+_S(f)=2^{-|T|}|{\cal F}_{|T}| .$$
Thus, in the language of traces, we are interested in
the effect of relaxing ``$\f$ shatters $Y$''
to require only that
${\cal F}_{|Y}$ contain a large fraction of $2^Y$.

The following arrow notation
(\eg, \cite {Bo,Fr,AMS19}) is convenient.  Write
$(N,n) \to (M,r)$ if every $\f\sub 2^{[n]}$
of size %
$N$   %
has a trace of size at least $M$ on some
$S \in \binom{[n]}{r}$; for example the Sauer--Shelah Theorem
says $(\Cc{n}{< r}+1,n) \to (2^r,r).$


\medskip
One might hope that Hamming balls would again give the
best examples in our relaxed setting, which would say, for example,
that for $m \le n$,
\begin {equation}
\label{e:tentc1}
(\Cc{n}{< r}+1,n) \to (\Cc{m}{< r}+1,m).
\end {equation}
But \eqref{e:tentc1}, which was first considered by
Bollob\'as and Radcliffe \cite{BR}
and would have implied both of the conjectures
disproved here,
was shown in \cite{BR} to be false
for fixed $r$ and (large) $m=n/2$.
(For $r=n/2$
and $m=n-1$, it fails for the original tribes example
discussed above.)



%
%
%
%
%
%
%
%
%
%
%
%
%
%

%


\subsubsection* {Two problems}


\begin{question}\label{prop4}

For fixed $\alpha,\delta >0$, what is
the largest $t\in (0,1/2)$ for which one can find
Boolean functions $f$ with $\mu(f) =t$ and
$I^+_S(f)<\ga$ for every $S\sub [n]$ of size
at most %
$(1/2-\gd)n$?
\end{question}

\nin
As far as we know $t> n^{-\gb}$ (with $\gb$ depending on $\ga,\gd$)
is possible. The influence of sets of half the variables is of special interest:

\begin{question}\label{prop5}

Given $\mu(f) =o(1)$ what can be said about 
the maximum of
the   %
$J_S^+(f)$'s for $|S|=n/2$?  %
What is the smallest $t$ 
such that for each $f$ with $\mu(f) =t$ there is some
$S$ of size $n/2$ with $J_S^+(f) \ge 1/2 ?$
(Here we assume that $n$ is even.)  %
\end{question}



\section{Boolean functions without influential coalitions}
\label{examples}


In each construction we consider, for suitable $k$ and $m$,
%
$f= \bigwedge_{i=1}^m C_i$, where the  %
%
$C_i$   %
%
%
%
are random %
disjunctions  %
of $k$ literals
using $k$ distinct variables (henceforth ``$k$-clauses'')
and show that $f$ is likely to have the desired properties. 
Every $C_i$ can be regarded as a list of specifications for 
the values of $k$ variables.
We use $g_i$ for the specification associated with $C_i$,
and write $C_i\sim x$ if some entry of $x$ agrees with $g_i$.
We say $C_i$ \emph{misses} $S\sub [n]$ if the indices of all
variables in $C_i$ lie in $[n]\sm S$.


Let $s= (1/2-\gd)n$.  %
We will always
use $S$ for an $s$-subset of $[n]$ and
(for such an $S$) set
$m_S=|\{i:\mbox{$C_i$ misses $S$}\}|.$
(Following common practice we omit irrelevant
floor and ceiling symbols,
pretending all large numbers are integers.
As in the case of $k,m$ and $s$, parameters not declared
to be constants are assumed to be functions of $n$.)
We use $\log$ for $\log_2$.


Both constructions will make use of the next two
observations, with \expref{Theorem}{prop1}
following immediately from these and
\expref{Theorem}{prop2} requiring a little more work.
\begin{lemma}\label{mSlemma}
If $k=o(\sqrt{n})$ and $(1/2+\gd)^km=\go(n)$ then
w.\,h.\,p.
\beq{mS}
%
%
%
%
%
%
%
m_S \sim (1/2+\gd)^km ~~ \text{ for all } S\in \Cc{[n]}{s}. %
%
\enq
\end{lemma}
\nin
(where, as usual, $a_n\sim b_n$ means $a_n/b_n\ra 1 $
and
\emph{with high probability} (w.\,h.\,p.) means
with probability tending to 1, both as $n\ra\infty$).


\mn
\emph{Proof.}
For a given $S$, $m_S$ has the binomial distribution
$B(m,p)$, with $p = \binom{n-s}{k}/\binom{n}{k}\sim(1/2+\gd)^k$
(using $k=o(\sqrt{n})$ for the ``$\sim $'').
Thus
$\mathbb E  m_S=mp$ and,
by ``Chernoff's Inequality''\footnote{Though now commonly called a ``Chernoff bound," this and related inequalities essentially go back to
 Sergei Natanovich  %
Bernstein in the early
part of the 20th century; see \cite{Wiki}.}
%
%
%
%
%
%
(\eg, \cite[Theorem 2.1]{JLR}),
\begin{equation}   \label{eq:chernoff}
\Pr(m_S\not\in ((1-\gz)mp,(1+\gz)mp)) < \exp [-\gO(\gz^2mp)],
\end{equation}
for $\gz\in (0,1)$.
Applying this with a $\gz$ which is both
$\go(\sqrt{n/(mp)})$ and $o(1)$ gives
$\Pr(m_S\not\sim mp )< 2^{-\go(n)}$,
and the union bound then gives \eqref{mS}.\qed


\medskip
The next lemma is stated to cover
both applications, though nothing so precise
is needed for \expref{Theorem}{prop1}.

\begin{lemma}\label{muflemma}
If there is a $\xi$ for which
\beq{xi1}
\exp[-\xi^2n] = o((1-2^{-k})^m)
\enq
and
\beq{xi2}
[(1+2\xi)/4]^k=o(1/m),
\enq
then
w.\,h.\,p.
\beq{muf}
\mu(f) \sim (1-2^{-k})^m.
\enq
\end{lemma}


\mn
\emph{Proof}.
This is a simple second moment method calculation
(similar to %
what is %
done in \cite{AL}, though described
differently there).


Recalling that $x,y$ always denote elements of $\{0,1\}^n$,
write $A_x$ for the event $\{f(x)=1\}$ and $\one_x$
for its indicator,
and set $X=\sum \one_x=2^n\mu(f) $.
Then
$\Pr(A_x)= (1-2^{-k})^m$ and
$\mathbb E  X = (1-2^{-k})^m2^n$; so we just need to show
$
\mathbb E  X^2\sim \mathbb E ^2 X
$ 
(equivalently, $\mathbb E  X^2<(1+o(1))\mathbb E ^2 X$),
since Chebyshev's Inequality then gives
$\Pr(|X-\mathbb E  X|> \gz\mathbb E  X) =o(1)$
for any fixed $\gz>0$.


We have
$$
\mathbb E  X^2 = \sum_x\sum_y\mathbb E  \one_x\one_y =\sum_x\Pr(A_x)\sum_y\Pr(A_y|A_x),
$$
so will be done if we show that for a fixed $x$,
$$ 
\sum_y\Pr(A_y|A_x) <(1+o(1)) (1-2^{-k})^m2^n.
$$ 
Since the sum is the same for all $x$, it is enough %
to prove this when $x=\uzero$.
Set $Z=\{y:|y|<(1/2-\xi)n\}$ and recall 
%
that by Chernoff's Inequality \eqref{eq:chernoff}, 
%
%
%
%
$|Z|<\exp[-2\xi^2n]2^n$.
It is thus enough to show that (for $x=\uzero$)
\beq{nicey}
y\not\in Z ~\Ra~\Pr(A_y|A_x)<(1+o(1)) (1-2^{-k})^m,
\enq
since then, using \eqref{xi1}, we have
$$
\sum_y\Pr(A_y|A_x) ~<~ |Z| + \sum_{y\not\in Z}\Pr(A_y|A_x)
~<~ (1+o(1)) (1-2^{-k})^m2^n.
$$


Now since $x=\uzero$, we have
%
$A_x =\{g_i\neq \uone ~~ \text{ for all } i\}$; %
%
%
so 
if, for a given $y\not\in Z$, we set
$\gb =\gb_y=\Pr(C_i\sim y|g_i\neq \uone)$
(a function of $|y|$),
then $\Pr(A_y|A_x) = \gb^m$.
Aiming for a bound on $\gb$, we have
%
\begin{align*}
1-2^{-k}&=\Pr(C_i\sim y) \\
&=
\Pr(g_i=\uone)\Pr(C_i\sim y|g_i=\uone)
+ \Pr(g_i\neq\uone)\Pr(C_i\sim y|g_i\neq \uone)\\
&= 2^{-k}\Pr(C_i\sim y|g_i=\uone) +(1-2^{-k})\gb
\end{align*}
and
$$\Pr(C_i\sim y|g_i=\uone) > 1 - (1-|y|/n)^k
\geq 1-(1/2+\xi)^k=:1-\nu
$$
(using the fact that if
$g_i=\uone$, then $g_i\not\sim y$
iff all indices of
variables in $C_i$ belong to $\{j:y_j=0\}$).
Combining, we have
$$
\gb < (1-2^{-k})^{-1}[1-2^{-k} - 2^{-k}(1-\nu)]
=(1-2^{-k})\left[1+\tfrac{2^{-k}(\nu -2^{-k})}{(1-2^{-k})^2}\right],
$$
which with \eqref{xi2} gives
$\gb^m < (1+o(1))(1-2^{-k})^m$ (which is \eqref{nicey}).\qed

\mn
\emph{Proof of \expref{Theorem}{prop1}}.    %
Notice that it is %
enough to prove this with $\mu(f) \sim \ga$
(rather than ``$=\ga$'');
for then,
since $f^{-1}(1)\sub g^{-1}(1)$ trivially implies
$J_S^+(f)\leq J_S^+(g)$ for all $S$, we can
choose
$\gb\in (\ga,1)$ and a $g$ with $\mathbb E (g)\sim \gb$
possessing the desired small influences, and
shrink $g^{-1}(1)$ to produce $f$.


Let $k=C\log n$, with $C=C_\gd$ chosen so that
$(1+2\gd)^k=\go(n)$ (\eg, $C=1/\gd$ does this),
and
$m = 2^k \ln (1/\ga)=n^C\ln (1/\ga)$.
Here all we use from \expref{Lemma}{mSlemma}
(whose hypotheses are satisfied for our choice of $k$ and $m$)
is the fact that
w.\,h.\,p. $m_S\neq 0$ for all $S$,
whence each $J_S^+(f)$ is at most
$1-2^{-k}=1-n^{-C}$.
On the other hand, by \expref{Lemma}{muflemma}
(with, for example, $\xi=0.1$), we have
$\mu(f) \sim \ga$ w.\,h.\,p.
So w.\,h.\,p. $f$ meets our requirements.\qed


\mn
%
%
\begin{proof}[Proof of \expref{Theorem}{prop2}].   %
Here, intending to recycle $n$, $m$ and $f$, we rename
these quantities
$\nn$, $\mm$ and $\ff$.
We may of course assume $\gd$ is fairly small.
Let (for example) $\xi = \gd/3$, fix $\eps$ with
$0<\eps < \xi^2$, and set
$k=(1+\gd)\log \nn$ and $\mm = \eps 2^k \nn$.
These values are easily seen to give the hypotheses of
\expref{Lemmas}{mSlemma} \expref{and}{muflemma}.
In particular, we can say that w.\,h.\,p.
the supports of the %
$C_i$  %
are chosen so \eqref{mS} holds
(note this says nothing about the values specified by
%
the $g_i$) %
and
\beq{mS'}
\mathbb E (\ff)\sim (1-2^{-k})^\mm\sim \eee^{-\eps \nn}.
\enq




Set $n=(1/2+\gd) \nn$.
Fix %
$S\in \binom{[\nn]}{s}$, %
set $m=m_S$,
and
let $f=f_S$ be the %
conjunction  %
of the $m$ %
clauses $C_i$%
---w.l.o.g. %
$C_1,\dots, C_m$%
---that miss $S$.
Thus $f$ is the %
conjunction of %
$m\sim (1/2+\gd)^k\mm =\eps (1+2\gd)^k\nn$
random $k$-clauses from a universe of $n$ variables.
\expref{Theorem}{prop2} (with $\ga=\gd$) thus follows from

\mn
\emph{Claim A.}
$\Pr(\mu(f)  > \exp[-n^{\gd}]) < o(2^{-\nn})$

\mn
(since then w.\,h.\,p. we have $\mathbb E (f_S)\leq \exp[-n^{\gd}]$ for every $S$).


\mn
\emph{Remarks.}
The actual bound in Claim A will be
$\exp[-\gO(m)]$,
so much smaller than $2^{-\nn}$.
Note that here it %
does not %
matter whether we take $\mu$ to be
our original measure (\ie, $\mu$ uniform on $\{0,1\}^\nn$)
or uniform measure on $Q:=\{0,1\}^n$;
but it is %
now more natural to think of
the latter---and we will do so in what follows---since
our original universe plays no further role in this discussion.
It may also be worth noting that, unlike in the proof of
\expref{Lemma}{muflemma}, the second moment method is not strong
enough to give the exponential bound in Claim A.




\mn
\emph{Claim B.}
If $X\sub Q$,
$\mu (X) =\gb> \exp[-o(n/\log^2n)]$ and $\gz =o(2^{-k})$,
then for a random $k$-clause $C$,
$$
\Pr(\mu (C\wedge X) >(1-\gz)\mu (X)) < 1/2
$$
(where $C\wedge X =\{x\in X: C\sim x\}$).


\mn
\emph{Remark.}  This is probably true for
$\gb$ greater than something like $\exp[-n/k]$.
The bound in the claim
is just what the proof gives, and is more than enough
for us since %
we are %
really interested in much larger $\gb$.


\medskip
To see that Claim B implies Claim A,
set $f_j=\bigwedge_{i=1}^jC_i$ and notice that $\mu(f) \geq \gb$
implies (for example)
\beq{imufi}
%
%
\left|\left\{i: \mathbb E (f_i) <\left(1-\tfrac{5\ln(1/\gb)}{m}\right)
\mathbb E (f_{i-1})\right\}\right| < m/5
%
\enq
(and, of course, $\mathbb E (f_i)\geq \gb$ for all $i$).
But if we take
$\gb =\exp[-n^{\gd}]$ then our choice of parameters gives
$$\gz:=5m^{-1}\ln(1/\gb) =o(2^{-k})$$
(using $m/\ln (1/\gb) =
\Theta(n^{(1+\gd)\log (1+2\gd)+1-\gd})$ and
$2^k = n^{1+\gd}$),
so Claim B bounds the probability of \eqref{imufi} by
%
\begin{equation*}
\Cc{m}{m/5} 2^{-4m/5} =o(2^{-\nn}).   \qedhere    %
\end{equation*}
%
\end{proof}   %

\mn
\emph{Proof of Claim B.}
Let $G$ be the bipartite graph on
$Q\cup W$, where $W$ is the set of $(n-k)$-dimensional
subcubes of $Q$ and, for $(x,D)\in Q\times W$,
we take $x\sim D$ if $x\in D$.
(So %
we have  %
gone to complements:  for a clause $C$
the corresponding subcube is $D=\{y:C\not\sim y\}$,
so $C\wedge X =X\sm D$.)

Assuming Claim B fails at $X$, fix $\SSS \sub W$ with $|\SSS |=|W|/2$ and
$$D\in \SSS
  ~\Ra ~\mu (D\cap X) < \gz\mu(X),$$
and set $\TTT =W\sm \SSS $.

%
Consider the experiment:  %
(i) choose $x$ uniformly from $X$;
(ii)  choose $D$ uniformly from the members of $W$ containing $x$;
(iii)  choose $y$ uniformly from $D$.


\mn
\emph{Claim C.}  $\Pr(y\in X)> (2-o(1))\gb$.


\mn
\begin{proof}   %
%
Since each triple $(x,D,y)$ with $x\in X$ and $x,y\in D$ is
produced by (i)-(iii) with probability
$|X|^{-1}\cdot 2^k|W|^{-1}\cdot 2^{k-n}$,
we just need to show that the number of such triples with
$y\in X$ is at least
$$
(2-o(1))\gb |X||W|2^n2^{-2k}  = (2-o(1))|X|^2|W|2^{-2k}.
$$
Writing $d$ for degree in $G$, we have
$$
\sum_{x\in X}d_\SSS(x) =\sum_{D\in \SSS }d_X(D) < |\SSS |\gz
 |X|,
$$
implying
%
\begin{align*}
\sum_{D\in \TTT }d_X(D)& = 
\sum_{x\in X}(d(x)-d_\SSS (x))\\
&> |X||W|2^{-k} - \gz |\SSS ||X| = (1-o(1))|X||W|2^{-k}.
\end{align*}   %
The number of %
triples $(x,D,y)$  %
as above is thus
%
\begin{align*}
\sum_{D\in W} d_X^2(D)
&\geq
\sum_{D\in \TTT } d_X^2(D)
\geq
\left(\sum_{D\in \TTT } d_X(D)\right)^2/|\TTT |\\   %
&> (1-o(1))|X|^2|W|^2|\TTT |^{-1}2^{-2k}
= (2-o(1))|X|^2|W|2^{-2k}\,.
 \qedhere                          
\end{align*}   %
%
\end{proof}   %

Let $T(x)$ be
the random element of $Q$ %
obtained %
from $x$
by choosing $K$ uniformly from $\binom{[n]}{k}$ and
randomly (uniformly, independently) 
resampling %
the %
$x_i$  %
%
%
%
for $i\not\in K$.    %
Then $y$ %
obtained  %
from $x$ by (ii) and (iii) above is just $T(x)$,
so
the next assertion contradicts Claim C, completing the proof of
Claim B (and \expref{Theorem}{prop2}).



\mn
\emph{Claim D.}
If $\mu(X) > \exp[-o(n/\log^2n)]$
and $x$ is uniform from $X$, then $\Pr(T(x)\in X) < (1+o(1))\mu(X)$.

\mn
\emph{Remark.}
If $X$ is a subcube of codimension $n/k$, say
$X=\{x:x \equiv 0 ~\mbox{on}~L\}$ with $|L |=n/k$,
then for any $x\in X$,
$$
\Pr(T(x)\in X) =\sum_t\Pr(|K\cap L |=t)2^{-(|L|-t)}
=\mu(X)\sum_t\Pr(|K\cap L |=t)2^t,
$$
and, since $|K\cap L|$ is essentially Poisson with mean 1,
the sum is
approximately
$ \eee^{-1}\sum_t2^t/t! = \eee$.
So Claim D fails for $\mu(X) = 2^{-n/k}$ and, as earlier,
it is %
natural to guess that it holds if $\mu(X)$ is much
%
greater  %
than this.

\mn
\emph{Proof of Claim D.}
Let $Q_r=\{y\in Q:|y|\leq r\}$.
The assumption on $\mu(X)$ implies that
$\mu(Q_{r-1})<\mu(X) \leq \mu(Q_r)$
for some $r >(1/2-o(1/k))n$, so Claim D follows from

\mn
\emph{Claim E.}
If $\vp =o(1/k)$ and $r > (1/2-\vp)n$, then for any $x\in Q$ and
$X\sub Q$ with
\beq{muX}
\mu(X) \leq \mu(Q_r),
\enq
we have

\beq{PrTx}
\Pr(T(x)\in X) < (1+o(1))\mu(Q_r).
\enq

\mn
\emph{Proof}.
We may assume $x=\underline{0}$, so that $\Pr(T(x)=y)$ is a decreasing
function of $|y|$.
We thus maximize $\Pr(T(x)\in X)$
subject to \eqref{muX} by taking
$
X=Q_r,
$
and \eqref{PrTx} is then a routine calculation
using
$$\mu(X) =\Pr({\rm Bin}(n,1/2)\leq r)$$
and
$$\Pr(T(x)\in X) =\Pr({\rm Bin}(n-k,1/2)\leq r)$$
(where ${\rm Bin}(\cdot,\cdot)$ denotes a binomially
distributed r.v.).\qed


\mn
{\bf Acknowledgment}
We would like to thank Roy Meshulam for helpful conversations.
Part of this work was carried out while the second author
was visiting Jerusalem
under BSF grant 2006066.
The first author
%
%
thanks the UC Berkeley %
Mathematics Department  %
for its hospitality.

%
